% !TeX program = pdflatex
%==============================================================================
%  Lista Teórica 08 --- Classificação e Classificadores Gaussianos
%  O gabarito sai de "Lista teorica 08 - gabarito.tex".
%==============================================================================
\providecommand{\gabaritoopt}{}
\documentclass[11pt,a4paper]{article}
\usepackage[\gabaritoopt]{../../recursos/latex/estilo-lista}

\begin{document}
\cabecalholista{08}{Classificação e Classificadores Gaussianos}

%------------------------------------------------------------------------------
\begin{exercicio}[o classificador de Bayes, em uma dimensão]
Seja $Y\in\{0,1\}$ e, condicionalmente a $Y$,
\[
  X\mid Y=0 \sim N(-1,\,1), \qquad X\mid Y=1 \sim N(+1,\,1).
\]
\begin{enumerate}[label=(\alph*),leftmargin=1.1cm]
  \item Com $\Prob(Y=1)=\Prob(Y=0)=1/2$, mostre que o classificador de Bayes
        decide por $Y=1$ quando $x>0$.
  \item Calcule o \textbf{erro de Bayes} nesse caso, em função de $\Phi$ (a
        distribuição acumulada da normal padrão), e dê o valor numérico.
        \emph{Dado:} $\Phi(-1)=0{,}1587$.
  \item Agora suponha $\Prob(Y=1)=0{,}1$. Mostre que o corte passa a ser
        $x^\ast=\tfrac12\ln 9 \approx 1{,}0986$, e calcule o novo erro de Bayes.
        \emph{Dados:} $\Phi(-2{,}0986)=0{,}0179$ e $\Phi(0{,}0986)=0{,}5393$.
  \item Compare o resultado de (c) com o erro do classificador que responde
        sempre $Y=0$. O que isso antecipa sobre a Aula~09?
\end{enumerate}
\end{exercicio}

\begin{solucao}
\textbf{(a)} O classificador de Bayes escolhe a classe de maior probabilidade a
posteriori, o que com duas classes equivale a decidir por $Y=1$ quando
$\pi_1 f_1(x) > \pi_0 f_0(x)$. Com $\pi_0=\pi_1$ as prioris cancelam e a condição
vira $f_1(x)>f_0(x)$, isto é
\[
  \exp\Big(-\tfrac{(x-1)^2}{2}\Big) > \exp\Big(-\tfrac{(x+1)^2}{2}\Big)
  \iff -(x-1)^2 > -(x+1)^2
  \iff (x+1)^2 > (x-1)^2 .
\]
Expandindo, $x^2+2x+1 > x^2-2x+1 \iff 4x>0 \iff x>0$. O corte é o ponto médio
entre as duas médias, como a simetria já sugeria.

\smallskip
\textbf{(b)} O erro é a probabilidade de errar, decomposta pelas duas classes:
\[
  \risco^\ast = \tfrac12\,\Prob(X>0\mid Y=0) + \tfrac12\,\Prob(X<0\mid Y=1).
\]
Sob $Y=0$, $X\sim N(-1,1)$, então $\Prob(X>0)=\Prob(Z>1)=1-\Phi(1)=\Phi(-1)$. Por
simetria, o segundo termo é igual. Logo
\[
  \risco^\ast = \Phi(-1) = \mathbf{0{,}1587}.
\]
Quase 16\% de erro é o \emph{piso}: nenhum classificador, com quaisquer dados,
faz melhor. É o análogo em classificação do $\sigma^2$ da Aula~01.

\smallskip
\textbf{(c)} Agora as prioris não cancelam. Decidimos por $Y=1$ quando
\[
  0{,}1\,\exp\Big(-\tfrac{(x-1)^2}{2}\Big) > 0{,}9\,\exp\Big(-\tfrac{(x+1)^2}{2}\Big).
\]
Tomando logaritmo,
\[
  \ln 0{,}1 - \tfrac{(x-1)^2}{2} > \ln 0{,}9 - \tfrac{(x+1)^2}{2}
  \iff \ln\tfrac{0{,}1}{0{,}9} > \tfrac{(x-1)^2-(x+1)^2}{2} = -2x,
\]
ou seja $-\ln 9 > -2x$, isto é $x > \tfrac12\ln 9 \approx 1{,}0986$.

O erro fica
\[
  \risco^\ast = 0{,}9\,\Prob(X>x^\ast\mid Y=0) + 0{,}1\,\Prob(X<x^\ast\mid Y=1)
   = 0{,}9\,\Phi(-2{,}0986) + 0{,}1\,\Phi(0{,}0986),
\]
usando $\Prob(X>x^\ast\mid Y=0)=1-\Phi(x^\ast+1)=\Phi(-x^\ast-1)$ e
$\Prob(X<x^\ast\mid Y=1)=\Phi(x^\ast-1)$. Numericamente,
\[
  \risco^\ast = 0{,}9\times 0{,}0179 + 0{,}1\times 0{,}5393
   = 0{,}0161 + 0{,}0539 = \mathbf{0{,}0701}.
\]

\smallskip
\textbf{(d)} O classificador que responde sempre $Y=0$ erra exatamente
$\Prob(Y=1)=\mathbf{0{,}10}$. O classificador de Bayes erra $0{,}0701$ --- melhor,
mas não muito: a diferença é de menos de 3 pontos percentuais.

E vale olhar de onde vem o erro de Bayes: das duas parcelas, $0{,}0539$ vêm de
classificar como $0$ observações que eram $1$. Ou seja, o classificador ótimo
**erra mais da metade dos positivos** ($0{,}5393$ deles!) --- e ainda assim é
ótimo, porque acertar mais positivos custaria muitos falsos positivos entre os
$90\%$ de negativos.

É exatamente a situação em que a acurácia deixa de ser informativa, e é o assunto
da Aula~09. Quando a prevalência é desbalanceada, um número de acurácia próximo
do trivial pode esconder tanto um classificador excelente quanto um inútil.
\end{solucao}

%------------------------------------------------------------------------------
\begin{exercicio}[por que a fronteira do LDA é uma reta]
No modelo gaussiano com covariância \textbf{comum},
$X\mid Y=k \sim N(\bm\mu_k,\,\bm\Sigma)$ com prioris $\pi_k$.
\begin{enumerate}[label=(\alph*),leftmargin=1.1cm]
  \item Escreva $\ln\big(\pi_k f_k(\x)\big)$ e mostre que, ao comparar duas
        classes, os termos que dependem de $\x$ de forma \emph{quadrática}
        se cancelam.
  \item Conclua que o classificador decide pela classe que maximiza a
        \emph{função discriminante linear}
        \[
          \delta_k(\x) = \x^\top\bm\Sigma^{-1}\bm\mu_k
            - \tfrac12\bm\mu_k^\top\bm\Sigma^{-1}\bm\mu_k + \ln\pi_k .
        \]
  \item Em que ponto exato do item (a) a hipótese $\bm\Sigma_0=\bm\Sigma_1$ é
        usada? O que sobra se ela for relaxada?
  \item No caso $d=1$, $\bm\Sigma=\sigma^2$, $\mu_0=-1$, $\mu_1=+1$ e
        $\pi_0=\pi_1$, recupere o corte do Exercício~1(a) a partir de
        $\delta_1(x)=\delta_0(x)$.
\end{enumerate}
\end{exercicio}

\begin{solucao}
\textbf{(a)} A densidade gaussiana multivariada dá
\[
  \ln\big(\pi_k f_k(\x)\big)
  = \ln\pi_k - \tfrac{d}{2}\ln(2\pi) - \tfrac12\ln\abs{\bm\Sigma}
    - \tfrac12(\x-\bm\mu_k)^\top\bm\Sigma^{-1}(\x-\bm\mu_k).
\]
Expandindo a forma quadrática:
\[
  (\x-\bm\mu_k)^\top\bm\Sigma^{-1}(\x-\bm\mu_k)
  = \underbrace{\x^\top\bm\Sigma^{-1}\x}_{\text{não depende de }k}
    - 2\,\x^\top\bm\Sigma^{-1}\bm\mu_k
    + \bm\mu_k^\top\bm\Sigma^{-1}\bm\mu_k .
\]
Ao comparar duas classes $k$ e $\ell$, subtraímos as duas expressões. Os termos
$-\tfrac d2\ln(2\pi)$, $-\tfrac12\ln\abs{\bm\Sigma}$ e
$-\tfrac12\x^\top\bm\Sigma^{-1}\x$ são \textbf{idênticos} nas duas e desaparecem
--- e é justamente o último deles que era quadrático em $\x$.

\smallskip
\textbf{(b)} O que sobra depois de descartar os termos comuns e multiplicar por
$-\tfrac12\cdot(-2)=1$ é
\[
  \delta_k(\x) = \x^\top\bm\Sigma^{-1}\bm\mu_k
    - \tfrac12\bm\mu_k^\top\bm\Sigma^{-1}\bm\mu_k + \ln\pi_k ,
\]
que é \textbf{afim em $\x$}: um produto interno com o vetor fixo
$\bm\Sigma^{-1}\bm\mu_k$, mais uma constante. A fronteira entre duas classes é o
conjunto $\{\x: \delta_k(\x)=\delta_\ell(\x)\}$, isto é
$\x^\top\bm\Sigma^{-1}(\bm\mu_k-\bm\mu_\ell) = \text{const}$ --- um
\textbf{hiperplano}.

\smallskip
\textbf{(c)} A hipótese é usada exatamente ao afirmar que
$\x^\top\bm\Sigma^{-1}\x$ e $\ln\abs{\bm\Sigma}$ \emph{não dependem de $k$}. Com
covariâncias distintas teríamos $\x^\top\bm\Sigma_k^{-1}\x$ e
$\ln\abs{\bm\Sigma_k}$, que dependem de $k$ e portanto \textbf{não se cancelam}.

Sobra então
\[
  \delta_k(\x) = -\tfrac12\x^\top\bm\Sigma_k^{-1}\x
    + \x^\top\bm\Sigma_k^{-1}\bm\mu_k
    - \tfrac12\bm\mu_k^\top\bm\Sigma_k^{-1}\bm\mu_k
    - \tfrac12\ln\abs{\bm\Sigma_k} + \ln\pi_k ,
\]
que é \emph{quadrática} em $\x$ --- é o QDA, e a fronteira vira uma cônica
(elipse, parábola ou hipérbole, conforme o caso). Em uma frase: \textbf{a
linearidade do LDA não vem da gaussianidade, vem da covariância comum}.

\smallskip
\textbf{(d)} Com $d=1$, $\bm\Sigma^{-1}=1/\sigma^2$ e $\pi_0=\pi_1$ (então
$\ln\pi_k$ cancela):
\[
  \delta_1(x)=\frac{x\cdot 1}{\sigma^2}-\frac{1}{2\sigma^2},
  \qquad
  \delta_0(x)=\frac{x\cdot(-1)}{\sigma^2}-\frac{1}{2\sigma^2}.
\]
Igualando: $x/\sigma^2 = -x/\sigma^2$, ou seja $2x/\sigma^2=0$, isto é
$\mathbf{x=0}$. É o corte do Exercício~1(a), e note que ele \emph{não depende de
$\sigma^2$} --- a variância muda o quanto se erra, não onde fica a fronteira.
\end{solucao}

%------------------------------------------------------------------------------
\begin{exercicio}[quantos parâmetros cada um estima]
Considere $K=2$ classes em $\R^d$.
\begin{enumerate}[label=(\alph*),leftmargin=1.1cm]
  \item Conte os parâmetros livres do LDA e do QDA. Considere: as médias das
        classes, a(s) matriz(es) de covariância --- lembrando que uma matriz
        simétrica $d\times d$ tem $d(d+1)/2$ entradas livres --- e as prioris.
  \item Preencha a tabela para $d=2$, $d=10$ e $d=50$.
  \item Mostre que a diferença QDA menos LDA vale $d(d+1)/2$ e comente a ordem de
        grandeza.
  \item A tabela abaixo traz a acurácia medida (média sobre 60 amostras, avaliada
        em 20 mil observações novas) do LDA e do QDA, nas mesmas duas dimensões.
        Em ambas o QDA passa a ganhar a partir de $n=50$ --- mas o que acontece
        \emph{antes} disso é muito diferente. Use a contagem de (b) para
        explicar, e diga por que a regra ``prefira LDA com poucos dados'' está
        incompleta.
\end{enumerate}

\begin{center}\small
\renewcommand{\arraystretch}{1.2}
\begin{tabular}{@{}llrrrrrr@{}}
\toprule
        &     & $n=20$ & $30$ & $50$ & $100$ & $300$ & $1000$ \\
\midrule
\multirow{2}{*}{$d=2$}  & LDA & $0{,}8602$ & $0{,}8681$ & $0{,}8724$ & $0{,}8775$ & $0{,}8785$ & $0{,}8793$ \\
                        & QDA & $0{,}8508$ & $0{,}8669$ & $0{,}8737$ & $0{,}8806$ & $0{,}8824$ & $0{,}8834$ \\
\midrule
\multirow{2}{*}{$d=10$} & LDA & $0{,}7053$ & $0{,}7422$ & $0{,}7773$ & $0{,}8013$ & $0{,}8175$ & $0{,}8241$ \\
                        & QDA & $0{,}5475$ & $0{,}7030$ & $0{,}7896$ & $0{,}8440$ & $0{,}8792$ & $0{,}8905$ \\
\bottomrule
\end{tabular}
\end{center}
\end{exercicio}

\begin{solucao}
\textbf{(a)} Com $K=2$:
\begin{itemize}[leftmargin=1.4cm]
  \item \textbf{médias}: dois vetores de $\R^d$, logo $2d$ parâmetros, nos dois
        métodos;
  \item \textbf{covariância}: o LDA estima \emph{uma} matriz comum,
        $d(d+1)/2$; o QDA estima \emph{duas}, $2\cdot d(d+1)/2 = d(d+1)$;
  \item \textbf{prioris}: $\pi_0$ e $\pi_1$ somam 1, então há $K-1=1$ parâmetro
        livre.
\end{itemize}
\[
  \#\text{LDA} = 2d + \tfrac{d(d+1)}{2} + 1,
  \qquad
  \#\text{QDA} = 2d + d(d+1) + 1 .
\]

\smallskip
\textbf{(b)}
\begin{center}\small
\renewcommand{\arraystretch}{1.3}
\begin{tabular}{@{}lrrr@{}}
\toprule
$d$ & $2$ & $10$ & $50$ \\
\midrule
LDA & $8$  & $76$  & $1\,376$ \\
QDA & $11$ & $131$ & $2\,651$ \\
\midrule
diferença & $3$ & $55$ & $1\,275$ \\
\bottomrule
\end{tabular}
\end{center}

\smallskip
\textbf{(c)} A diferença é
$\big(2d+d(d+1)+1\big)-\big(2d+\tfrac{d(d+1)}{2}+1\big)=\tfrac{d(d+1)}{2}$: a
segunda matriz de covariância, exatamente. Ela cresce como $d^2/2$, ou seja
\textbf{quadraticamente}. Em $d=2$ são 3 parâmetros a mais --- nada. Em $d=50$ são
$1\,275$ a mais, quase o dobro do modelo inteiro do LDA.

\smallskip
\textbf{(d)} O ponto de virada é o mesmo nas duas dimensões --- $n=50$ --- e
justamente por isso ele é a parte \emph{menos} informativa da tabela. O que muda
é o \textbf{tamanho} da vantagem de cada lado:

\begin{center}\small
\renewcommand{\arraystretch}{1.2}
\begin{tabular}{@{}lcc@{}}
\toprule
                                    & $d=2$ & $d=10$ \\
\midrule
LDA $-$ QDA em $n=20$ (antes da virada)   & $+0{,}0094$ & $\mathbf{+0{,}1578}$ \\
QDA $-$ LDA em $n=1000$ (depois)          & $+0{,}0041$ & $\mathbf{+0{,}0664}$ \\
custo da matriz extra, $d(d+1)/2$         & $3$ & $55$ \\
\bottomrule
\end{tabular}
\end{center}

Em $d=2$ os dois métodos são praticamente equivalentes em toda a tabela: a maior
diferença é de menos de um ponto percentual. A matriz de covariância extra custa
$3$ parâmetros, e $3$ parâmetros a mais ou a menos não mudam nada, nem para bem
nem para mal.

Em $d=10$ os dois erram feio, em direções opostas. Com $n=20$, o QDA precisa
estimar $131$ parâmetros a partir de $20$ observações --- as matrizes de
covariância mal são invertíveis --- e despenca para $0{,}5475$, quase o acaso,
$16$ pontos abaixo do LDA. Com $n=1000$ ele reproduz a estrutura verdadeira e
ganha $6{,}6$ pontos.

Portanto a regra ``prefira LDA com poucos dados'' está incompleta em dois
sentidos. Primeiro, \textbf{o que decide não é $n$ sozinho, e sim $n$ contra
$d(d+1)/2$}, que é o custo da matriz extra. Segundo, e mais útil na prática:
\textbf{em dimensão baixa a escolha quase não importa}; ela só passa a importar
quando $d$ é grande, e aí importa muito, nos dois sentidos. Gastar validação
cruzada decidindo entre LDA e QDA em $d=2$ é desperdício; em $d=50$ é obrigatório.

É a mesma lógica do balanço viés--variância da Aula~01: o QDA é o modelo mais
flexível, e modelos flexíveis só compensam quando há dados para sustentá-los.
\end{solucao}

%------------------------------------------------------------------------------
\begin{exercicio}*[naive Bayes gaussiano é QDA com as mãos amarradas]
O \texttt{GaussianNB} supõe que, \emph{dentro de cada classe}, as covariáveis são
independentes: $\bm\Sigma_k$ é \textbf{diagonal}, com entradas
$\sigma_{k1}^2,\dots,\sigma_{kd}^2$.
\begin{enumerate}[label=(\alph*),leftmargin=1.1cm]
  \item Quantos parâmetros ele estima, com $K=2$? Acrescente a coluna à tabela
        do Exercício~3.
  \item Mostre que a função discriminante do naive Bayes gaussiano é
        \[
          \delta_k(\x) = \ln\pi_k
            - \tfrac12\sum_{j=1}^{d}\left[\ln\sigma_{kj}^2
              + \frac{(x_j-\mu_{kj})^2}{\sigma_{kj}^2}\right],
        \]
        e diga se a fronteira é linear ou quadrática.
  \item A suposição de independência é quase sempre \emph{falsa}. Dê uma razão
        pela qual o método ainda assim costuma classificar bem.
  \item Descreva uma situação em que a suposição custa caro, isto é, em que a
        correlação entre covariáveis carrega informação essencial sobre a classe.
\end{enumerate}
\end{exercicio}

\begin{solucao}
\textbf{(a)} São $2d$ médias, $2d$ variâncias (uma por covariável e por classe) e
$1$ priori: $\#\text{NB}=4d+1$. Comparando:
\begin{center}\small
\renewcommand{\arraystretch}{1.3}
\begin{tabular}{@{}lrrr@{}}
\toprule
$d$ & $2$ & $10$ & $50$ \\
\midrule
naive Bayes ($4d+1$) & $9$  & $41$  & $201$ \\
LDA                  & $8$  & $76$  & $1\,376$ \\
QDA                  & $11$ & $131$ & $2\,651$ \\
\bottomrule
\end{tabular}
\end{center}
Note a inversão: em $d=2$ o naive Bayes custa \emph{mais} que o LDA (9 contra 8),
porque estima duas variâncias por covariável em vez de uma covariância comum. Em
$d=50$ ele custa $201$ contra $1\,376$ --- quase sete vezes menos. O crescimento
é \textbf{linear} em $d$, contra o quadrático dos outros dois, e é isso que o
torna viável em dimensão alta (a Aula~E3 o usa com dezenas de milhares de
colunas).

\smallskip
\textbf{(b)} Com $\bm\Sigma_k$ diagonal, a densidade conjunta fatora:
$f_k(\x)=\prod_{j=1}^d f_{kj}(x_j)$, com
$f_{kj}$ a densidade de $N(\mu_{kj},\sigma_{kj}^2)$. Então
\[
  \ln\big(\pi_k f_k(\x)\big)
   = \ln\pi_k + \sum_{j=1}^{d}\left[
       -\tfrac12\ln(2\pi) - \tfrac12\ln\sigma_{kj}^2
       - \frac{(x_j-\mu_{kj})^2}{2\sigma_{kj}^2}\right],
\]
e descartando $-\tfrac d2\ln(2\pi)$, que é comum a todas as classes, chega-se à
expressão do enunciado.

A fronteira é \textbf{quadrática}: o termo $(x_j-\mu_{kj})^2/\sigma_{kj}^2$
depende de $k$ pelo denominador, então o termo em $x_j^2$ não cancela entre
classes. O naive Bayes gaussiano é, portanto, um \textbf{QDA restrito a
covariâncias diagonais} --- e não uma versão simplificada do LDA. (Se as
variâncias fossem forçadas a ser iguais entre as classes, aí sim ele viraria um
LDA diagonal, com fronteira linear.)

\smallskip
\textbf{(c)} Porque para \emph{classificar} basta acertar qual $\delta_k$ é o
maior, e não estimar bem as probabilidades. Erros na modelagem da dependência
tendem a deslocar todos os $\delta_k$ na mesma direção, preservando a ordem entre
eles. Some-se a isso o argumento de variância do item (a): com $4d+1$ parâmetros
em vez de $2\,651$, as estimativas são muito mais estáveis, e o viés introduzido
pela suposição falsa é compensado pela variância que se economiza.

O preço aparece na \textbf{calibração}: as probabilidades do naive Bayes são
notoriamente ruins, tipicamente empurradas para 0 ou 1, porque ele multiplica
evidências correlacionadas como se fossem independentes e conta a mesma
informação várias vezes. É exatamente a medição da Aula~09 --- AUC praticamente
empatada com a logística ($0{,}776$ contra $0{,}774$) e Brier muito pior
($0{,}1523$ contra $0{,}0653$). \emph{Ordenar bem não é estimar bem.}

\smallskip
\textbf{(d)} Quando a informação está na \textbf{correlação}, e não nas médias.
Um exemplo em $d=2$: suponha que nas duas classes cada covariável tenha média
zero e variância um, mas que na classe 0 elas sejam positivamente correlacionadas
($\rho=+0{,}9$) e na classe 1, negativamente ($\rho=-0{,}9$). As nuvens são duas
elipses cruzadas, perfeitamente distinguíveis --- mas as médias e as variâncias
marginais são \emph{idênticas} nas duas classes.

O naive Bayes vê exatamente médias e variâncias marginais, então para ele as duas
classes são indistinguíveis. Simulando $4\,000$ observações por classe e medindo
a acurácia por validação cruzada de 5 dobras:
\begin{center}\small
\begin{tabular}{@{}lc@{}}
\toprule
\texttt{GaussianNB} & $0{,}5104$ \\
LDA                 & $0{,}4949$ \\
QDA                 & $\mathbf{0{,}8575}$ \\
\bottomrule
\end{tabular}
\end{center}
O naive Bayes e o LDA acertam o mesmo que jogar uma moeda; o QDA, que estima as
duas matrizes de covariância inteiras, chega a $86\%$. É o caso em que a
suposição de independência não custa precisão --- custa \emph{tudo}.

(O LDA falha pelo mesmo motivo, e é instrutivo: ele estima uma covariância
\emph{comum}, e a média das duas matrizes tem correlação $\approx 0$. A
informação está justamente na \emph{diferença} entre elas.)
\end{solucao}

\paraalem

\end{document}
